\section*{Capítulo \ref{ch:g10:synthesis-yield} --- A síntese: rendimento e pureza}

\begin{solution}{exo:g10:synthesis-yield:1}
$\eta = 4.2 / 5.0 = 0.84$, um \omterm{def:g10:synthesis-yield:yield}{rendimento} de \qty{84}{\percent}.
\end{solution}

\begin{solution}{exo:g10:synthesis-yield:2}
Pesagem dos \omterm{def:g7:chemical-reactions:reactant}{reagentes}; \omterm{def:g10:synthesis-yield:reflux}{aquecimento sob refluxo}; \omterm{def:g10:synthesis-yield:vacuum-filtration}{filtração a vácuo} do
\omterm{def:g10:synthesis-yield:synthesis}{produto bruto}; \omterm{def:g10:synthesis-yield:recrystallisation}{recristalização}; medição do ponto de fusão.
\end{solution}

\begin{solution}{exo:g10:synthesis-yield:3}
Entrando por baixo, a água enche a camisa inteira antes de sair por
cima, e assim o condensador é resfriado em todo o seu comprimento. O
alto fica aberto para que a montagem não fique fechada: aquecer uma
montagem fechada acumularia pressão e poderia fazê-la estourar.
\end{solution}

\begin{solution}{exo:g10:synthesis-yield:4}
O produto deve ser muito solúvel no \omterm{def:g6:solutions-and-solubility:solute}{solvente} quente e pouco solúvel no
\omterm{def:g6:solutions-and-solubility:solute}{solvente} frio. As impurezas, presentes em pequenas quantidades,
continuam dissolvidas no \omterm{def:g6:solutions-and-solubility:solute}{solvente} frio e acabam no \omterm{def:g3:separating-mixtures:filter}{filtrado}.
\end{solution}

\begin{solution}{exo:g10:synthesis-yield:5}
Não: ela funde abaixo de \qty{135}{\celsius} e ao longo de uma faixa de
vários graus, sinal de um sólido impuro. Ela deve ser recristalizada,
secada, e o seu ponto de fusão medido de novo.
\end{solution}

\begin{solution}{exo:g10:synthesis-yield:6}
$n = 2.00 / 138.0 = \qty{0.0145}{mol}$ de ácido salicílico, o \omterm{def:g10:reaction-progress-table:limiting}{reagente
limitante}; $m_{\max} = 0.0145 \times 180.0 = \qty{2.61}{g}$;
$\eta = 2.03 / 2.61 = 0.78$, isto é, \qty{78}{\percent}.
\end{solution}

\begin{solution}{exo:g10:synthesis-yield:7}
A massa pesada inclui água, por isso é maior que a massa de aspirina: o
\omterm{def:g10:synthesis-yield:yield}{rendimento} calculado fica alto demais, aqui impossivelmente acima de
\qty{100}{\percent}. O \omterm{def:g10:synthesis-yield:synthesis}{produto bruto} seco ainda contém impurezas (ácido
salicílico que não reagiu, subprodutos), que também acrescentam massa:
o \omterm{def:g10:synthesis-yield:yield}{rendimento} dele também fica alto demais em relação ao \omterm{def:g10:synthesis-yield:yield}{rendimento}
verdadeiro em aspirina pura.
\end{solution}

\begin{solution}{exo:g10:synthesis-yield:8}
Ela é muito mais rápida e deixa os cristais muito mais secos, já que a
sucção retira a maior parte do líquido; os cristais também são fáceis de
lavar no funil e de raspar do papel de filtro plano.
\end{solution}

\begin{solution}{exo:g10:synthesis-yield:9}
Produto $3.1/5.0 = 0.62$; aspirina $3.1/5.0 = 0.62$; ácido salicílico
$2.2/5.0 = 0.44$. O produto se desloca como a aspirina e não mostra
mancha de ácido salicílico: é aspirina, sem ácido salicílico detectado.
\end{solution}

\begin{solution}{exo:g10:synthesis-yield:10}
\omterm{def:g6:solutions-and-solubility:solute}{Solvente} A: muito solúvel a quente, pouco solúvel a frio; o \omterm{def:g6:solutions-and-solubility:solute}{solvente} B
mantém a maior parte do produto dissolvida mesmo a frio. Com A,
\qty{3.0}{g} precisam de cerca de $3.0 / 150 = \qty{0.020}{L}$, isto é,
\qty{20}{mL} de \omterm{def:g6:solutions-and-solubility:solute}{solvente} quente; depois de frio, no máximo
$2 \times 0.020 = \qty{0.04}{g}$ continua dissolvido.
\end{solution}

\begin{solution}{exo:g10:synthesis-yield:11}
Um líquido em ebulição precisa de espaço acima dele: a espuma e os
respingos não podem chegar ao condensador. As pedrinhas de ebulição dão
ao vapor pequenos lugares onde formar bolhas, para que o líquido ferva
de modo regular em vez de em jatos repentinos.
\end{solution}

\begin{solution}{exo:g10:synthesis-yield:12}
$0.80 \times 0.70 = 0.56$: \qty{56}{\percent}. Três etapas de
\qty{90}{\percent}: $0.90^3 = 0.73$, \qty{73}{\percent}. Cada etapa perde
produto (e tempo e dinheiro), e as perdas se multiplicam: quanto menos
etapas, mais produto sobra no fim.
\end{solution}

\begin{solution}{exo:g10:synthesis-yield:13}
No máximo $3.3 \times 0.050 = \qty{0.17}{g}$ de aspirina. A água gelada
dissolve ainda menos aspirina, já que a sua \omterm{def:g6:solutions-and-solubility:solubility}{solubilidade} diminui com a
temperatura.
\end{solution}

\begin{solution}{exo:g10:synthesis-yield:14}
\begin{enumerate}
  \item $n_0(\text{ácido salicílico}) = 2.76 / 138.0 = \qty{0.0200}{mol}$,
        anidrido \qty{0.0500}{mol}: o ácido salicílico é limitante,
        $x_{\max} = \qty{0.0200}{mol}$, e sobram $0.0500 - 0.0200 =
        \qty{0.0300}{mol}$ de anidrido.
  \item A \omterm{def:g10:synthesis-yield:synthesis}{síntese} dá \qty{0.0200}{mol} de ácido etanoico, a destruição
        do excesso $2 \times 0.0300 = \qty{0.0600}{mol}$:
        \qty{0.0800}{mol} ao todo.
\end{enumerate}
\end{solution}

\begin{solution}{exo:g10:synthesis-yield:15}
$n(\text{aspirina}) = \num{1.00e6} / 180.0 = \qty{5.56e3}{mol}$. Com um
\omterm{def:g10:synthesis-yield:yield}{rendimento} de \qty{85}{\percent}, ácido salicílico necessário:
$\num{5.56e3} / 0.85 = \qty{6.54e3}{mol}$, isto é,
$\num{6.54e3} \times 138.0 = \qty{9.0e5}{g}$, cerca de \qty{0.90}{t}.
\end{solution}

\begin{solution}{pb:g10:synthesis-yield:1}
\textbf{1.} C: $7 + 4 = 11$ e $9 + 2 = 11$; H: $6 + 6 = 12$ e
$8 + 4 = 12$; O: $3 + 3 = 6$ e $4 + 2 = 6$. Balanceada.

\textbf{2.} Ácido salicílico $7 \times 12.0 + 6 \times 1.0 + 3 \times 16.0
= \qty{138.0}{g/mol}$; anidrido $4 \times 12.0 + 6 \times 1.0 +
3 \times 16.0 = \qty{102.0}{g/mol}$; aspirina $9 \times 12.0 + 8 \times
1.0 + 4 \times 16.0 = \qty{180.0}{g/mol}$; ácido etanoico $2 \times 12.0
+ 4 \times 1.0 + 2 \times 16.0 = \qty{60.0}{g/mol}$.

\textbf{3.} \omterm{def:g7:chemical-reactions:reactant}{Reagentes}: o ácido salicílico e o anidrido etanoico; produto
desejado: a aspirina; subproduto: o ácido etanoico.

\textbf{4.} O aquecimento torna a reação mais rápida; sob refluxo, os
vapores são condensados e devolvidos, e assim nenhum \omterm{def:g7:chemical-reactions:reactant}{reagente} ou produto
se perde.

\textbf{5.} $n = 3.0 / 138.0 = \qty{0.0217}{mol}$.

\textbf{6.} $m = 6.0 \times 1.08 = \qty{6.48}{g}$;
$n = 6.48 / 102.0 = \qty{0.0635}{mol}$.

\textbf{7.} Ácido salicílico: $0.0217 - x$; anidrido $0.0635 - x$;
aspirina e ácido etanoico: $x$. O ácido salicílico se esgota primeiro:
$x_{\max} = \qty{0.0217}{mol}$, o ácido salicílico é limitante.

\textbf{8.} $m_{\max} = 0.0217 \times 180.0 = \qty{3.91}{g}$.

\textbf{9.} Para que todo o ácido salicílico seja transformado em
aspirina: o que sobrasse ficaria \omterm{def:g2:mixing-and-dissolving:mix}{misturado} à aspirina e seria difícil de
retirar, enquanto o anidrido em excesso é destruído pela água
acrescentada no fim, e o ácido etanoico que ele dá é eliminado na
lavagem.

\textbf{10.} Ela inclui água: \qty{3.6}{g} dariam um falso
$3.6 / 3.91 = \qty{92}{\percent}$.

\textbf{11.} $3.3 / 3.91 = 0.84$: \qty{84}{\percent} de \omterm{def:g10:synthesis-yield:synthesis}{produto bruto}
seco, que não é puro.

\textbf{12.} As impurezas são retiradas, e uma parte da aspirina continua
dissolvida no \omterm{def:g6:solutions-and-solubility:solute}{solvente} frio e nas águas de lavagem.

\textbf{13.} $3.3 \times 0.030 = \qty{0.10}{g}$ no máximo.

\textbf{14.} $\eta = 2.7 / 3.91 = 0.69$: \qty{69}{\percent}.

\textbf{15.} Um ponto de fusão nítido no valor da aspirina pura: os
cristais são aspirina, e pura.

\textbf{16.} Fusão abaixo de \qty{135}{\celsius} e ao longo de uma faixa
larga: o \omterm{def:g10:synthesis-yield:synthesis}{produto bruto} era impuro.

\textbf{17.} O \omterm{def:g10:synthesis-yield:synthesis}{produto bruto} continha aspirina e ácido salicílico que não
reagiu; a \omterm{def:g10:synthesis-yield:recrystallisation}{recristalização} retirou o ácido salicílico.

\textbf{18.} $R_f = 3.1 / 5.0 = 0.62$.

\textbf{19.} Mais ou menos o mesmo em cada uma: de \qty{3.91}{g}
possíveis a \qty{3.3}{g} de \omterm{def:g10:synthesis-yield:synthesis}{produto bruto}, e de \qty{3.3}{g} a
\qty{2.7}{g} na \omterm{def:g10:synthesis-yield:recrystallisation}{recristalização}, cerca de \qty{0.6}{g} em cada uma. Como
o \omterm{def:g10:synthesis-yield:synthesis}{produto bruto} não era aspirina pura, a primeira parte do trabalho na
verdade perdeu um pouco mais de \qty{0.6}{g} de aspirina.

\textbf{20.} O \omterm{def:g10:synthesis-yield:yield}{rendimento} da \omterm{def:g10:synthesis-yield:synthesis}{síntese} é de cerca de \qty{69}{\percent}.
\end{solution}
